\documentclass[11pt]{article}
\usepackage[T1]{fontenc}
\usepackage{textcomp}
\usepackage[margin=0.75in]{geometry}
\usepackage{amsmath,amssymb,amsthm}
\usepackage{array,booktabs}
\usepackage{hyperref}
\setlength{\parindent}{0pt}
\newtheorem{theorem}{Theorem}
\theoremstyle{definition}
\newtheorem{definition}{Definition}
\title{Flow Proof Artifact: Nat-plus-zero-right}
\author{Generated by \texttt{flow doc proof}}
\date{Flow Proof Book}
\begin{document}
\maketitle
n + 0 = n, derived by induction.\\[0.5em]
\textbf{Source.} peano/induction — Gries \& Schneider, Ch. 3\\[0.5em]
\bigskip
\noindent{\large \textbf{Derived fact 1.} \textit{Adding zero on the right gives the same number} \hfill the natural numbers \cdot addition \cdot zero is the right identity}\par
\smallskip\noindent\textit{Coordinate: the natural numbers \cdot addition \cdot zero is the right identity}\par
\smallskip\noindent\textit{Source: peano/induction — https://en.wikipedia.org/wiki/Mathematical\_induction}\par
\smallskip\noindent\textit{Built on: adding zero on the left does not change the number, adding one more on the right bumps the sum by one}\par
\medskip
\noindent\textbf{Goal.} Adding zero on the right gives the same number.  (12 + 0 = 12)\par
\noindent$\displaystyle \forall n \in \mathbb{N}\quad n + 0 = n$\par
\medskip
\noindent
\begin{tabular}{@{} >{\raggedright\arraybackslash}p{0.06\textwidth} >{\raggedright\arraybackslash}p{0.47\textwidth} >{\raggedleft\arraybackslash}p{0.41\textwidth} @{}}
\textbf{\#} & \textbf{Proof} & \textbf{Mathematics} \\
\midrule
\textbf{1.} & We prove that zero is the right identity for addition on the natural numbers. & \\[0.45em]
\textbf{2.} & We proceed by induction on n: first the base case, then the inductive step. & \\[0.45em]
\textbf{3.} & Consider the base case in which n is zero. & \\[0.45em]
\textbf{4.} & We invoke the definitional clause governing addition on the natural numbers: adding zero on the left does not change the number (instantiated for 0). & $\displaystyle 0 + 0 = 0$ \\[0.45em]
\textbf{5.} & From step 3 and step 4, we can deduce that 0 plus 0 equals 0. This establishes the base case (see step 3 and step 4). Hence proven. & $\displaystyle 0 + 0 = 0$ \\[0.45em]
\textbf{6.} & For the inductive step, suppose the claim holds for all smaller values. & \\[0.45em]
\textbf{7.} & Under the supposition in step 6, let k denote the predecessor of n. & \\[0.45em]
\textbf{8.} & Under the supposition in step 6, we cross the inductive boundary: assume the claim holds for k (the induction hypothesis). & $\displaystyle k + 0 = k$ \\[0.45em]
\textbf{9.} & We invoke the definitional clause governing addition on the natural numbers: adding one more on the right bumps the sum by one (instantiated for k, 0). & $\displaystyle k + \mathrm{succ}(0) = \mathrm{succ}(k + 0)$ \\[0.45em]
\textbf{10.} & From step 6, step 7, step 8, and step 9, this implies n plus 0 equals n. Hence proven. & $\displaystyle n + 0 = n$ \\[0.45em]
\bottomrule
\end{tabular}
\label{thm:Nat-addition-zero_is_the_right_identity}
\par\medskip\hrule
\end{document}
